\section{Introduction}

This paper develops and evaluates \textbf{TAReL} (Temporal Adversarial
Reconstruction of High-Frequency Latent Spaces in Financial Markets), a
pipeline for studying limit-order-book (LOB) microstructure through the
geometry of its recurrence plots and through the seed-to-seed
reproducibility of generative models trained on those plots. Three
lineages of work intersect with the proposal developed here:
\begin{itemize}\setlength{\itemsep}{1pt}
\item Generative modeling of market microstructure dynamics.
\item Recurrence-plot generation and analysis of financial time series.
\item Reproducibility analysis for adversarial training in finance.
\end{itemize}
Each lineage is briefly contrasted with our setting; the goal is not
to displace any of these lines but to compose them into a unified
approach.

The premise the paper is built on is the following: \emph{the
recurrence structure of an order book --- the geometric pattern of
when a market briefly returns to a state it has visited before ---
encodes something invariant about that market's microstructure.} Not
the prices themselves, not the trades, not the volumes, but the shape
of how the order book moves through state space over time. If true,
that shape is a fingerprint. Two markets with similar microstructure
regimes should produce similar recurrence patterns; two markets with
dissimilar microstructure should produce visibly different ones. The
conjecture is unfalsifiable in the abstract but sharply testable
through the pipeline we develop here.
